

\begin{frame}{Parlons transaction et réplication}

Une \alert{transaction} est une séquence d'opérations de \alert{lecture} ou d'\alert{écriture}, 
se terminant par \sqlkw{commit} ou \sqlkw{rollback}.


\vfill

\begin{redbullet}

  \item   Le \texttt{commit} est une instruction qui \alert{valide} toutes les mises à jour.

  \item Le \texttt{rollback} est une instruction qui \alert{annule} toutes les mises à jour.

\end{redbullet}
\vfill

\begin{remblock}{Essentiel}
Les opérations d'une transaction sont solidaires: elles sont toutes validées, ou pas du tout
(\alert{atomicité}).
\end{remblock}
\end{frame}


\begin{frame}{Propriétés des transactions (en relationnel)}

Un système relationnel contrôle la concurrence et
garantit un ensemble de propriétés rassemblées sous l'acronyme ACID.

\vfill

\begin{redbullet}

  \item  \alert{A = Atomicité}. Une transaction est validée complètement ou pas du tout.
  
  \item  \alert{C = Cohérence}. Une transaction mène d'un état cohérent à un autre état cohérent.
  
  \item \alert{I = Isolation}. Une transaction s'exécute comme si elle était seule.
  
   \item \alert{D = Durabilité}. Quand le commit s'exécute, ses résultats sont définitifs.
\end{redbullet}

\begin{noteblock}
Les algorithmes sont \alert{très} complexes, et impactent les performances.
\end{noteblock}
\end{frame}



\begin{frame}{Et en NoSQL}

Les transactions sont indispensables en relationnel à cause de la multiplicité des
écritures.
\vfill

En NoSQL, c'est différent

\begin{redbullet}

  \item  On peut souvent lire / écrire un seul document\,: pas besoin de transaction
  
  \item On a souvent des applications analytiques de type \textit{Write once, read many (WORM)}. 
  
  \item Et \alert{surtout} les systèmes NoSQL sont conçus pour être
      distribués, et ça devient très très compliqué 
  
\end{redbullet}

\vfill

Les systèmes NoSQL ne proposent pas (sauf exception) des transactions
ACID, mais un protocole simplifié, dit BASE.

\end{frame}

\section{La réplication, pourquoi?}

\begin{frame}
\frametitle{La réplication}

Outil indispensable, universel pour la robustesse des systèmes distribués.

\begin{redbullet}
  \item \alert{Tolérance aux pannes}
  
     Vous cherchez un document sur $S_1$, qui est en panne\,? On le trouvera sur $S_2$
     
 \item \alert{Distribution des lectures}
 
    Répartissons les lectures sur $S_1, S_2, \ldots, S_n$ pour satisfaire les millions de requêtes de nos clients.
    
 \item \alert{Distribution des écritures\,?}
 
    Oui, mais attention, il faut \alert{réconcilier} les données ensuite.
    
 \end{redbullet}

Le bon niveau de réplication? \alert{Trois copies pour une sécurité totale; deux au minimum}.

\end{frame}


\begin{frame}
\frametitle{Méthode générale de réplication}

Une \alert{application} (le client) demande au \alert{système}  $S_1$ \alert{l'écriture}
d'un document.

\begin{redbullet}
  \item $S_1$ écrit le document sur le disque\,; 
  \item \alert{$S_1$ transmet la demande d'écriture à un ou plusieurs autres serveurs $S_2, \ldots, S_n$}, créant
     des \alert{replicas} ou \alert{copies}\,;
  \item $S_1$ ``rend la main'' au client (acquittement)
\end{redbullet}

Essentiel, réplication synchrone/asynchrone\,?

\begin{redbullet}
  \item \alert{synchrone}: $S_1$ \alert{attend} confirmation de $S_2, \ldots, S_n$ avant de rendre la main\,; 
 \item \alert{asynchrone}: $S_1$ rend la main \alert{sans attendre}. 
 \end{redbullet}

\begin{remblock}{Conséquences}
Une réplication \alert{asynchrone} est plus rapide, mais  favorise les \alert{incohérences}.
\end{remblock}


\end{frame}

\begin{frame}
\frametitle{Réplication avec écritures synchrones}

Le client est acquitté quand \alert{tous} les serveurs ont effectué l'écriture.

\vfill

\figSlideWithSize{replication-synchrone}{10cm}

\vfill

Sécurité totale; cohérence forte; \alert{plus lent}.

\end{frame}

\begin{frame}
\frametitle{Réplication avec écritures asynchrones}

Le client est acquitté quand \alert{un} des serveurs a effectué l'écriture. Les autres
écritures se font indépendamment.

\vfill

\figSlideWithSize{replication-asynchrone}{10cm}

\vfill

Sécurité partielle; cohérence faible; \alert{efficace}.

\end{frame}

\section{Cohérence des données}


\begin{frame}
\frametitle{Parlons cohérence}

La cohérence est la capacité d’un système de gestion de données à refléter fidèlement les opérations d’une application.
 
\begin{redbullet}
  \item toute opération (validée) est immédiatement visible et permanente. 
  \item si je fais une écriture de $d$ suivie d’une lecture, je dois constater les modifications effectuées; 
  \item si je refais une lecture un peu plus tard, ces modifications doivent toujours être présentes.
\end{redbullet}

\alert{Cohérence forte = } une des fonctionnalités marquantes des systèmes relationnels.

\begin{remblock}{Dans les systèmes NoSQL}
\alert{La cohérence forte est sacrifiée au profit des performances.} 
\end{remblock}
\end{frame}


\begin{frame}
\frametitle{Les transactions de type BASE}

BASE est un modèle de transaction adapté aux systèmes distribués.

\vfill

\begin{redbullet}
  \item \alert{Basically Available} -- Pas de verrouillage, on privilégie
       la disponibilité des données
       
  \item \alert{Soft State} -- on accepte des états intermédiaires décalés par rapport 
      aux mises à jour en cours.
  
  \item \alert{Eventually Consistent} (Cohérence à terme) ) -- le système
     garantit que les incohérences ne sont que transitoires. 
\end{redbullet}
\vfill

Une propriété souvent cité (théorème CAP) dit qu'on ne peut pas tout avoir...
\vfill

\begin{noteblock}{C'est important\,?}
BASE est un modèle transactionnel faible et problématique pour des 
application à mise à jour intensive. Sans impact pour des applications
analytiques
\end{noteblock}

\end{frame}



\begin{frame}
\frametitle{À retenir}

\begin{redbullet}
  \item L'absence de transactions ACID, comme l'absence de jointure, est
une faiblesse importante des systèmes NoSQL pour les applications
généralistes.

  \item Ces systèmes priorisent la capacité à traiter de grands
   volumes de données dans un contexte distribué.

  \item Pour les applications analytiques traitant des données
     figées en lecture, pas de problème.
  
\end{redbullet}

\end{frame}

\end{document}


